\documentclass[a4paper,12pt]{ctexart}
\usepackage{geometry}
\usepackage{amsmath}
\usepackage{amssymb}
\usepackage{enumitem}
\usepackage{booktabs}
\usepackage{tikz}
\usetikzlibrary{arrows.meta}

% 页面设置
\geometry{left=2.5cm, right=2.5cm, top=2.5cm, bottom=2.5cm}
\linespread{1.25}

% 全局首行缩进
\setlength{\parindent}{2em}
\setlength{\parskip}{0.3em}

% 列表设置
\setlist{
    leftmargin=2em,
    itemindent=0em,
    labelindent=0em,
    topsep=0.3em,
    itemsep=0.1em
}

% 矢量
\newcommand{\vct}[1]{\boldsymbol{#1}}
\newcommand{\uve}[1]{\boldsymbol{\hat{#1}}}

\title{格物助航计划\\质点运动学基础知识点提纲}
\author{北航工科物理资料整理组}
\date{\today}

\begin{document}

\maketitle

\vspace{1em}

\tableofcontents
\newpage

\section{第一章：质点运动学}

\subsection{章节主线与学习目标}
本章的核心任务是建立描述机械运动的定量语言，不涉及运动变化的原因。本部分（1.1-1.5）按照"描述前提-理想模型-数学工具-基本物理量"的逻辑展开：
\[
\text{时空观} \rightarrow \text{质点模型} \rightarrow \text{参考系与坐标系} \rightarrow \text{位置描述} \rightarrow \text{运动快慢描述}
\]

通过本部分学习，读者应具备以下核心能力：
\begin{enumerate}[label=\arabic*.]
\item 能够独立建立描述运动的完整坐标系，明确所有约定条件；
\item 能够准确判断任意物理场景下质点模型的适用性；
\item 能够熟练写出直角坐标系下质点的运动方程，并推导其轨道方程；
\item 能够严格区分矢量与标量物理量，正确计算位移、速度和速率。
\end{enumerate}

\subsection{时间与空间}
\textbf{重点学习}
\begin{enumerate}
\item 机械运动的本质是物体的空间位置随时间的变化，时间和空间是描述运动的两个基本维度。
\item 经典力学的时空观：时间和空间是绝对的、均匀的、各向同性的。时间均匀流逝，与参考系无关；空间是平坦的欧几里得空间，长度测量与参考系无关。
\item 参考系：为了确定物体的位置和描述其运动而选作标准的物体或物体系。运动的描述具有相对性，同一物体在不同参考系中的运动形式可能完全不同。
\item 坐标系：固结在参考系上的数学标架，由原点和一组正交单位矢量组成，用于将物体的位置量化为数值坐标。坐标系是参考系的数学抽象，同一参考系可以选择不同的坐标系。
\item 国际单位制（SI）：力学的三个基本物理量及其单位为：长度（米，m）、时间（秒，s）、质量（千克，kg）。所有力学量都可以由这三个基本量导出，量纲分析是检验物理公式正确性的重要方法。
\end{enumerate}

\textbf{常见误区}
\begin{itemize}
\item 认为参考系必须是静止的：参考系可以任意选择，通常选择使问题简化的参考系；
\item 混淆参考系和坐标系：参考系是实物，坐标系是数学工具，二者不能等同；
\item 忽略量纲检查：所有物理公式两边的量纲必须一致，量纲错误的公式一定是错误的。
\end{itemize}

\textbf{基础量纲分析入门}

量纲是物理量的基本属性，用于表示物理量与基本物理量之间的导出关系。在力学中，国际单位制的三个基本量纲为：长度$[L]$、时间$[T]$、质量$[M]$。任意力学量$A$的量纲均可表示为幂次组合形式：
\[
[A] = L^\alpha T^\beta M^\gamma
\]
其中$\alpha,\beta,\gamma$为量纲指数。

\textbf{量纲分析的核心作用}：
\begin{enumerate}[label=(\arabic*)]
\item 快速检验物理公式的正确性（量纲不一致的公式必然错误）；
\item 推导未知物理量的量纲，进而确定其国际单位；
\item 辅助建立物理量之间的定性关系。
\end{enumerate}

\textbf{重要提示}：量纲正确是公式成立的必要条件，而非充分条件。

\textbf{运动学量纲分析例题}
某同学推导匀变速直线运动的位移公式时得到了以下三个表达式，其中$x$为位移，$v_0$为初速度，$a$为加速度，$t$为运动时间。请用量纲分析判断哪个表达式一定是错误的：
\begin{enumerate}[label=(\Alph*)]
\item $x = v_0 t + \frac{1}{2} a t^2$
\item $x = v_0 t + a t$
\item $x = \frac{1}{2} a t^2$
\end{enumerate}

\textbf{解析}：
分别计算等式两边各项的量纲：
\begin{itemize}
\item 位移$x$的量纲：$[x] = L$
\item 速度$v_0$的量纲：$[v_0] = L T^{-1}$，因此$[v_0 t] = L T^{-1} \cdot T = L$
\item 加速度$a$的量纲：$[a] = L T^{-2}$，因此$[a t^2] = L T^{-2} \cdot T^2 = L$，$[a t] = L T^{-2} \cdot T = L T^{-1}$
\end{itemize}
选项B中，$v_0 t$的量纲为$L$，而$a t$的量纲为$L T^{-1}$，两项量纲不一致，不能进行代数相加，因此该表达式一定是错误的。

量纲分析是贯穿整个物理学的通用方法，其适用范围远不止力学领域。在电磁学、热学、光学、原子物理等所有物理分支中，所有物理量的量纲都可以由国际单位制（SI）规定的七个基本量纲导出。任意物理量$A$的通用量纲表达式为：
\[
[A] = L^\alpha T^\beta M^\gamma I^\delta \Theta^\varepsilon N^\zeta J^\eta
\]
其中$I$为电流的量纲，$\Theta$为热力学温度的量纲，$N$为物质的量的量纲，$J$为发光强度的量纲，$\alpha,\beta,\gamma,\delta,\varepsilon,\zeta,\eta$为对应的量纲指数。

\begin{table}[htbp]
  \centering
  \caption{国际单位制（SI）基本量纲与基本单位（2019年修订版）}
  \begin{tabular}{cccc}
    \toprule
    基本量纲 & 物理量名称 & 基本单位名称 & 单位符号 \\
    \midrule
    $L$ & 长度 & 米 & m \\
    $T$ & 时间 & 秒 & s \\
    $M$ & 质量 & 千克 & kg \\
    $I$ & 电流 & 安培 & A \\
    $\Theta$ & 热力学温度 & 开尔文 & K \\
    $N$ & 物质的量 & 摩尔 & mol \\
    $J$ & 发光强度 & 坎德拉 & cd \\
    \bottomrule
  \end{tabular}
  \label{tab:si_basic_units}
\end{table}

注：国际单位制中还包含弧度（rad）和球面度（sr）两个辅助单位，用于表示平面角和立体角，其量纲为1（无量纲）。

\textbf{理论力学衔接}

经典力学的时空观是整个牛顿力学的基础。在理论力学中，我们将进一步引入"构型空间"的概念，物体的所有可能位置构成一个多维空间，每个点对应物体的一个构型。

\subsection{物体的点模型：质点}
\textbf{重点学习}
\begin{enumerate}[label=\arabic*.]
\item 质点是力学中最基本的理想模型，定义为具有一定质量但几何尺寸可以忽略不计的物体。
\item 理想模型的意义：抓住问题的主要矛盾，忽略次要因素，使复杂的物理问题得以简化。
\item 质点模型的适用条件：当物体的大小和形状对所研究问题的影响可以忽略不计时，即可将物体视为质点。具体包括两种情况：
    \begin{enumerate}[label=(\alph*)]
    \item 物体的线度远小于研究问题中涉及的距离；
    \item 物体只做平动，不发生转动或形变。
    \end{enumerate}
\item 一个物体能否视为质点，完全取决于所研究的问题，与物体本身的大小无关。
\end{enumerate}

\textbf{常见误区}
\begin{itemize}
\item 认为小物体一定能视为质点，大物体一定不能视为质点：例如，研究原子的内部结构时，原子不能视为质点；研究地球绕太阳的公转时，地球可以视为质点；
\item 认为质点就是几何点：质点具有质量，而几何点没有质量；
\item 不考虑研究问题，笼统地说某个物体是质点。
\end{itemize}

\textbf{理论力学衔接}


质点模型是单自由度系统的最简单情况。在理论力学中，我们将研究更复杂的模型，如质点系（多个质点组成的系统）和刚体（形状和大小不变的物体）。刚体可以看作是质点系的一种特例，即刚体可以视作质点之间无相对位移的质点系。刚体的引入是理论力学学习的非常重要的内容，也是由理论迈向实际的关键一环。

刚体是由无数个质点组成、且任意两质点间距离保持不变的理想模型。对于自由刚体，除了需要3个坐标确定其质心的平动位置外，还需要3个独立的转角坐标（通常称为欧拉角）确定其在空间中的取向，因此自由刚体总共具有6个自由度：3个平动自由度 + 3个转动自由度。这也解释了质点模型的本质：当物体的转动和形变对所研究问题的影响可以忽略时，我们实际上是将6个自由度的刚体简化为仅具有3个平动自由度的质点，从而极大地简化了问题。

刚体自由度的概念在工程领域有极其广泛的应用：例如航空航天中卫星的姿态控制，本质上就是通过推力器和飞轮精确控制卫星的3个转动自由度（俯仰、偏航、滚转）；机械工程中的工业机器人，每个关节通常对应一个自由度，通过多个关节的组合实现末端执行器的任意位姿；而在研究车辆的平面运动时，我们通常只考虑3个自由度（质心的x、y坐标和车身的转角），并通过轮胎的约束条件进一步简化问题。在理论力学中，我们还将进一步学习约束对自由度的影响，以及如何使用广义坐标来描述具有任意自由度的力学系统。

\subsection{位置矢量与轨道方程}
\textbf{重点学习}
\begin{enumerate}[label=\arabic*.]
\item 位置矢量（位矢）：从坐标系原点指向质点所在位置的矢量，记为$\vct{r}$。位矢同时给出了质点到原点的距离和相对于原点的方位。
\item 直角坐标系中的位矢表示：
\[
\vct{r}(t) = x(t)\uve{i} + y(t)\uve{j} + z(t)\uve{k}
\]
其中$\uve{i},\uve{j},\uve{k}$分别为x、y、z轴正方向的单位矢量，$x(t),y(t),z(t)$为质点在三个坐标轴上的坐标，都是时间$t$的函数。
\item 运动方程：描述质点位置随时间变化规律的方程，其矢量形式为$\vct{r}=\vct{r}(t)$，分量形式为$x=x(t),y=y(t),z=z(t)$。运动方程包含了质点运动的全部信息。
\item 轨道方程：描述质点在空间中运动轨迹的几何方程。由运动方程的分量形式消去时间参数$t$，即可得到轨道方程$f(x,y,z)=0$。
\end{enumerate}

\textbf{常用坐标系及其位矢表示}
坐标系是描述质点位置的数学工具，选择合适的坐标系可以极大简化问题的求解。在运动学中，最常用的三种坐标系为笛卡尔坐标系、平面极坐标系和自然坐标系。

\subsubsection{笛卡尔坐标系（直角坐标系）}
笛卡尔坐标系是最基础、最通用的坐标系，由三条相互垂直的坐标轴组成。
\begin{itemize}
\item \textbf{定义}：以空间中某点$O$为原点，建立三条相互垂直的坐标轴$x$、$y$、$z$，分别对应三个正交单位矢量$\uve{i}$、$\uve{j}$、$\uve{k}$。
\item \textbf{适用情况}：绝大多数直线运动、直角坐标系下的曲线运动（如抛体运动）、以及没有明显对称性的一般运动。
\item \textbf{位矢表示}：
\[
\vct{r} = x\uve{i} + y\uve{j} + z\uve{k}
\]
\end{itemize}

\begin{figure}[htbp]
  \centering
  \begin{tikzpicture}[scale=1.2, >=Stealth]
    \draw[->] (0,0,0) -- (3,0,0) node[right] {$x$};
    \draw[->] (0,0,0) -- (0,3,0) node[above] {$y$};
    \draw[->] (0,0,0) -- (0,0,3) node[below left] {$z$};
    \draw[->, thick, blue] (0,0,0) -- (2,2,1) node[above right] {$\vct{r}$};
    \draw[dashed] (2,2,1) -- (2,2,0) -- (2,0,0);
    \draw[dashed] (2,2,0) -- (0,2,0);
    \node at (2.2,0,0) {$x$};
    \node at (0,2.2,0) {$y$};
    \node at (0,0,1.2) {$z$};
    \node at (1.5,0,0) {$\uve{i}$};
    \node at (0,1.5,0) {$\uve{j}$};
    \node at (0,0,1.5) {$\uve{k}$};
  \end{tikzpicture}
  \caption{笛卡尔坐标系与位矢}
\end{figure}

\subsubsection{平面极坐标系}
平面极坐标系适用于描述平面内具有径向对称性的运动。
\begin{itemize}
\item \textbf{定义}：以平面内某点$O$为原点，建立水平极轴。质点的位置由极径$r$（原点到质点的距离）和极角$\theta$（极轴到位矢的夹角）确定，对应两个正交单位矢量$\uve{e}_r$（径向，沿位矢 outward）和$\uve{e}_\theta$（角向，垂直于$\uve{e}_r$逆时针方向）。
\item \textbf{适用情况}：径向运动、圆周运动、中心力问题（如行星绕太阳的运动、卫星轨道）。
\item \textbf{位矢表示}：
\[
\vct{r} = r\uve{e}_r
\]
注意：极坐标中的单位矢量$\uve{e}_r$和$\uve{e}_\theta$是随质点位置变化的，这是与笛卡尔坐标系的本质区别。
\end{itemize}

\begin{figure}[htbp]
  \centering
  \begin{tikzpicture}[scale=1.5, >=Stealth]
    \draw[->] (-1,0) -- (3,0) node[right] {};
    \node at (3.1,0) {极轴};
    \node at (-0.2,-0.2) {$O$};
    \draw[->, thick, blue] (0,0) -- (45:2) node[above right] {$\vct{r} = r\uve{e}_r$};
    \node at (45:1) {$r$};
    \draw (0.5,0) arc(0:45:0.5) node[midway, above right] {$\theta$};
    \draw[->, thick, red] (45:2) -- (45:2.7) node[above right] {$\uve{e}_r$};
    \draw[->, thick, green] (45:2) -- (135:2.7) node[above left] {$\uve{e}_\theta$};
  \end{tikzpicture}
  \caption{平面极坐标系与位矢}
\end{figure}

\subsubsection{自然坐标系}
自然坐标系是沿质点运动轨迹建立的坐标系，特别适合描述已知轨迹的曲线运动，采用圆形轨迹演示最为直观。
\begin{itemize}
\item \textbf{定义}：在质点的运动轨迹上任选一点$O$为原点，规定轨迹的正方向，用弧长$s$表示质点的位置。在质点所在位置建立两个正交单位矢量：切向单位矢量$\uve{\tau}$（沿轨迹切线指向运动正方向）和法向单位矢量$\uve{n}$（沿轨迹法线指向凹侧，对于圆周运动即指向圆心）。
\item \textbf{适用情况}：已知轨迹的曲线运动（如圆周运动、摆线运动、车辆沿道路的运动）。
\item \textbf{位矢表示}：
\[
\vct{r} = \vct{r}(s)
\]
其中$s$为质点到原点$O$的弧长，是时间$t$的函数$s=s(t)$。
\end{itemize}

\begin{figure}[htbp]
  \centering
  \begin{tikzpicture}[scale=1.2, >=Stealth]
    \draw[thick] (0,0) circle (2);
    \node at (0,2.2) {轨迹正方向};
    \draw[->] (120:2) arc(120:90:2);
    \fill (2,0) circle (2pt) node[below right] {$O$};
    \fill (45:2) circle (2pt) node[above right] {$P$};
    \draw[->] (2,0) arc(0:45:2) node[midway, above right] {$s$};
    \draw[->, thick, red] (45:2) -- (135:2.7) node[above left] {$\uve{\tau}$};
    \draw[->, thick, green] (45:2) -- (45:1.3) node[below left] {$\uve{n}$};
  \end{tikzpicture}
  \caption{自然坐标系与位矢（圆形轨迹）}
\end{figure}

\textbf{常见误区}
\begin{itemize}
\item 混淆运动方程和轨道方程：运动方程是时间的函数，描述质点在任意时刻的位置；轨道方程是几何关系，只描述质点运动的轨迹形状，不包含时间信息；
\item 消去时间参数时忽略参数的取值范围，导致轨道方程的定义域错误；
\item 认为位矢的大小就是质点的位移：位矢是从原点指向质点的矢量，位移是质点在一段时间内位置的变化量。
\end{itemize}

\textbf{理论力学衔接}

在理论力学中，我们将引入广义坐标的概念。广义坐标是能够唯一确定系统构型的一组独立变量，不一定是直角坐标。例如，对于在平面上运动的质点，可以用极坐标$(r,\theta)$作为广义坐标。

\subsection{速度矢量}

\subsubsection{速度矢量的基础定义}
速度是描述质点运动快慢和方向的物理量，是矢量。为了准确理解速度的概念，我们首先明确以下几个基本物理量：

\textbf{位移}：质点在一段时间内位置的变化量，是矢量。设质点在$t$时刻的位矢为$\vct{r}(t)$，在$t+\Delta t$时刻的位矢为$\vct{r}(t+\Delta t)$，则在$\Delta t$时间内的位移为：
\[
\Delta\vct{r} = \vct{r}(t+\Delta t) - \vct{r}(t)
\]
位移的大小是质点初末位置之间的直线距离，方向由初位置指向末位置。

\textbf{路程}：质点在一段时间内实际运动轨迹的长度，是标量。一般情况下，位移的大小$|\Delta\vct{r}|$不等于路程$\Delta s$，只有当质点做单向直线运动时，二者才相等。

\textbf{平均速度}：位移与发生这段位移所用时间的比值，是矢量：
\[
\overline{\vct{v}} = \frac{\Delta\vct{r}}{\Delta t}
\]
平均速度的方向与位移方向相同，它只能粗略地描述质点在一段时间内的平均运动状态。

\textbf{瞬时速度（简称速度）}：平均速度在$\Delta t \to 0$时的极限，即位矢对时间的一阶导数：
\[
\vct{v} = \lim_{\Delta t \to 0} \frac{\Delta\vct{r}}{\Delta t} = \frac{d\vct{r}}{dt}
\]
速度是矢量，其方向沿质点运动轨迹的切线方向，指向运动的前方。

\textbf{速率}：速度的大小，是标量：
\[
v = |\vct{v}| = \left| \frac{d\vct{r}}{dt} \right| = \frac{ds}{dt}
\]
其中$s$为质点运动的路程。速率等于路程对时间的变化率，它只描述质点运动的快慢，不涉及方向。

\subsubsection{不同坐标系下的速度矢量}
速度矢量的表达式与坐标系的选择密切相关。基于前一节介绍的三种常用坐标系，我们分别推导速度的分量形式。

\textbf{1. 笛卡尔坐标系（直角坐标系）}
在笛卡尔坐标系中，单位矢量$\uve{i},\uve{j},\uve{k}$是恒定不变的，因此速度的推导最为简单。将位矢$\vct{r}=x\uve{i}+y\uve{j}+z\uve{k}$对时间求导：
\[
\vct{v} = \frac{d\vct{r}}{dt} = \frac{dx}{dt}\uve{i} + \frac{dy}{dt}\uve{j} + \frac{dz}{dt}\uve{k} = v_x\uve{i} + v_y\uve{j} + v_z\uve{k}
\]
其中$v_x=\frac{dx}{dt},v_y=\frac{dy}{dt},v_z=\frac{dz}{dt}$分别为速度在三个坐标轴上的分量。速度的大小为：
\[
v = \sqrt{v_x^2 + v_y^2 + v_z^2}
\]

\textbf{2. 平面极坐标系}
在平面极坐标系中，位矢表示为$\vct{r}=r\uve{e}_r$。与笛卡尔坐标系不同，极坐标的单位矢量$\uve{e}_r$和$\uve{e}_\theta$是随质点位置变化的，因此求导时必须考虑单位矢量本身的变化：
\[
\vct{v} = \frac{d\vct{r}}{dt} = \frac{dr}{dt}\uve{e}_r + r\frac{d\uve{e}_r}{dt}
\]
通过几何分析可以证明，径向单位矢量的导数为$\frac{d\uve{e}_r}{dt}=\frac{d\theta}{dt}\uve{e}_\theta$，代入上式得到极坐标系下的速度表达式：
\[
\vct{v} = \dot{r}\uve{e}_r + r\dot{\theta}\uve{e}_\theta = v_r\uve{e}_r + v_\theta\uve{e}_\theta
\]
其中$v_r=\dot{r}$称为径向速度，描述质点到原点距离的变化率；$v_\theta=r\dot{\theta}$称为角向速度（或横向速度），描述质点方位角的变化率。

\textbf{3. 自然坐标系}
在自然坐标系中，质点的位置由弧长$s$表示，位矢为$\vct{r}=\vct{r}(s)$。根据复合函数求导法则：
\[
\vct{v} = \frac{d\vct{r}}{dt} = \frac{d\vct{r}}{ds} \cdot \frac{ds}{dt}
\]
其中$\frac{d\vct{r}}{ds}$的大小为1，方向沿轨迹切线方向，即切向单位矢量$\uve{\tau}$。因此自然坐标系下的速度表达式为：
\[
\vct{v} = v\uve{\tau} = \frac{ds}{dt}\uve{\tau}
\]
自然坐标系下的速度只有切向分量，这是其最显著的特点。

\subsubsection{矢量求导的物理意义}
矢量求导是运动学中最核心的数学工具，它与标量求导有着本质的区别：
- 标量的导数只描述其大小的变化率
- 矢量的导数同时描述其大小的变化率和方向的变化率

对于任意矢量$\vct{A}(t)$，其导数可以分解为两部分：
\[
\frac{d\vct{A}}{dt} = \frac{dA}{dt}\uve{e}_A + A\frac{d\uve{e}_A}{dt}
\]
其中第一项$\frac{dA}{dt}\uve{e}_A$描述矢量大小的变化，第二项$A\frac{d\uve{e}_A}{dt}$描述矢量方向的变化。

在运动学中，这一区别至关重要：
- 当质点做直线运动时，速度方向不变，只有大小变化，因此加速度只改变速度大小
- 当质点做曲线运动时，速度方向不断变化，即使速度大小不变，也存在加速度（如匀速圆周运动）

矢量求导的思想贯穿整个力学课程，是后续学习加速度、角动量等物理量的基础。

\subsubsection{典型例题：匀速圆周运动的速度分析}
为了直观说明速度的矢量性，我们以匀速圆周运动为例进行分析。设质点在半径为$R$的圆周上以恒定的速率$v$运动，$t=0$时质点位于x轴正方向。

\textbf{方法一：笛卡尔坐标系}
质点的运动方程为：
\[
x = R\cos\omega t, \quad y = R\sin\omega t
\]
其中$\omega=\frac{v}{R}$为角速度。对时间求导得到速度分量：
\[
v_x = \frac{dx}{dt} = -R\omega\sin\omega t, \quad v_y = \frac{dy}{dt} = R\omega\cos\omega t
\]
速度的大小为：
\[
v = \sqrt{v_x^2 + v_y^2} = R\omega = \text{常量}
\]
速度的方向由方向角$\alpha$确定：
\[
\tan\alpha = \frac{v_y}{v_x} = -\cot\omega t = \tan\left(\omega t + \frac{\pi}{2}\right)
\]
可见速度方向始终与位矢方向垂直，即沿圆周切线方向。

\textbf{方法二：极坐标系}
在极坐标系中，$r=R=\text{常量}$，$\theta=\omega t$。代入极坐标速度公式：
\[
v_r = \dot{r} = 0, \quad v_\theta = r\dot{\theta} = R\omega = v
\]
因此速度只有角向分量，大小为$v$，方向沿圆周切线方向，与笛卡尔坐标系结果一致。

\textbf{结论}：匀速圆周运动的速度大小保持不变，但方向不断变化，因此它是一种变速运动，必然存在加速度。这一结论充分体现了速度的矢量性。

\textbf{常见误区}
\begin{itemize}
\item 混淆位移和路程：位移是矢量，路程是标量；
\item 混淆平均速度和平均速率：平均速率是路程与时间的比值，是标量；平均速度的大小是位移大小与时间的比值，一般不等于平均速率；
\item 认为速度的方向总是与位矢的方向相同：速度方向沿轨迹切线方向，与位矢方向无关；
\item 认为速度的大小不变时，速度就不变：速度是矢量，大小和方向只要有一个发生变化，速度就发生变化；
\item 忽略极坐标和自然坐标系中单位矢量的变化，求导时只对分量求导而不对单位矢量求导。
\end{itemize}

\textbf{理论力学衔接}

速度是描述质点运动状态的基本物理量之一。在理论力学中，我们将看到，质点的运动状态由其位置和速度共同决定。哈密顿力学则进一步引入动量作为基本变量，建立了更普适的力学体系。对于刚体的运动，我们还将引入角速度矢量来描述其转动状态。

\end{document}